\documentclass[10pt,a4paper]{amsart}
\usepackage{amsmath,amssymb,amsthm,graphicx}
% Retained minimal preamble for the Stacks Project chapter more-algebra.tex (More on Algebra).
% The chapter file pinned-original-source.tex carries only \input{preamble} on its first line, so no
% preamble text lives inside the chapter itself.  The Stacks Project uses one single shared preamble for
% every chapter; the snapshot stages it identically as sources/stacks-chow/preamble.tex and
% sources/stacks-etale-cohomology/preamble.tex.  The block below is copied mechanically, line for line,
% from sources/stacks-chow/preamble.tex lines 14-18 (the Xy-pic setup) and 154-217 (the native theorem
% declarations together with the mathematical macros).  Nothing was added, removed or redefined; the
% wrapper class and the amsmath/amssymb/amsthm/graphicx base packages are supplied by the assembler.
% Chapter: more-algebra.tex -- More on Algebra.
% For commutative diagrams we use Xy-pic
\usepackage[all]{xy}

% We use 2cell for 2-commutative diagrams.
\xyoption{2cell}

% Theorem environments.
%
\theoremstyle{plain}
\newtheorem{theorem}[subsection]{Theorem}
\newtheorem{proposition}[subsection]{Proposition}
\newtheorem{lemma}[subsection]{Lemma}

\theoremstyle{definition}
\newtheorem{definition}[subsection]{Definition}
\newtheorem{example}[subsection]{Example}
\newtheorem{exercise}[subsection]{Exercise}
\newtheorem{situation}[subsection]{Situation}

\theoremstyle{remark}
\newtheorem{remark}[subsection]{Remark}
\newtheorem{remarks}[subsection]{Remarks}

\numberwithin{equation}{subsection}

% Macros
%
\def\lim{\mathop{\mathrm{lim}}\nolimits}
\def\colim{\mathop{\mathrm{colim}}\nolimits}
\def\Spec{\mathop{\mathrm{Spec}}}
\def\Hom{\mathop{\mathrm{Hom}}\nolimits}
\def\Ext{\mathop{\mathrm{Ext}}\nolimits}
\def\SheafHom{\mathop{\mathcal{H}\!\mathit{om}}\nolimits}
\def\SheafExt{\mathop{\mathcal{E}\!\mathit{xt}}\nolimits}
\def\Sch{\mathit{Sch}}
\def\Mor{\mathop{\mathrm{Mor}}\nolimits}
\def\Ob{\mathop{\mathrm{Ob}}\nolimits}
\def\Sh{\mathop{\mathit{Sh}}\nolimits}
\def\NL{\mathop{N\!L}\nolimits}
\def\CH{\mathop{\mathrm{CH}}\nolimits}
\def\proetale{{pro\text{-}\acute{e}tale}}
\def\etale{{\acute{e}tale}}
\def\QCoh{\mathit{QCoh}}
\def\Ker{\mathop{\mathrm{Ker}}}
\def\Im{\mathop{\mathrm{Im}}}
\def\Coker{\mathop{\mathrm{Coker}}}
\def\Coim{\mathop{\mathrm{Coim}}}

% Boxtimes
%
\DeclareMathSymbol{\boxtimes}{\mathbin}{AMSa}{"02}

%
% Macros for moduli stacks/spaces
%
\def\QCohstack{\mathcal{QC}\!\mathit{oh}}
\def\Cohstack{\mathcal{C}\!\mathit{oh}}
\def\Spacesstack{\mathcal{S}\!\mathit{paces}}
\def\Quotfunctor{\mathrm{Quot}}
\def\Hilbfunctor{\mathrm{Hilb}}
\def\Curvesstack{\mathcal{C}\!\mathit{urves}}
\def\Polarizedstack{\mathcal{P}\!\mathit{olarized}}
\def\Complexesstack{\mathcal{C}\!\mathit{omplexes}}
% \Pic is the operator that assigns to X its picard group, usage \Pic(X)
% \Picardstack_{X/B} denotes the Picard stack of X over B
% \Picardfunctor_{X/B} denotes the Picard functor of X over B
\def\Pic{\mathop{\mathrm{Pic}}\nolimits}
\def\Picardstack{\mathcal{P}\!\mathit{ic}}
\def\Picardfunctor{\mathrm{Pic}}
\def\Deformationcategory{\mathcal{D}\!\mathit{ef}}

% --- appended 2026-10-12 (Chief): the project own definitions of the reference/slogan/history
% display environments, copied mechanically from the shared Stacks preamble lines 8-12
% (\usepackage{verbatim} plus the three \newenvironment lines, which the chapter preamble keeps
% verbatim from the project preamble). Nothing added, removed or redefined: these environments
% carry no mathematical content and the project itself typesets them as comments. Without these
% three lines a counted statement whose own text uses them fails with Environment slogan undefined.
\usepackage{verbatim}
\newenvironment{reference}{\comment}{\endcomment}
\newenvironment{slogan}{\comment}{\endcomment}
\newenvironment{history}{\comment}{\endcomment}
\usepackage{cleveref}
\crefname{originallemma}{original lemma}{original lemmas}
\Crefname{originallemma}{Original lemma}{Original lemmas}
\makeatletter
\newlabel{algebra-lemma-NAK}{{00DV}{source}{}{sourcepointer-1}{}}
\newlabel{algebra-lemma-NAK@cref}{{[originallemma][0][]00DV}{[1][0][]source}{}{}{}}
\makeatother
\crefname{originallemma}{original lemma}{original lemmas}
\Crefname{originallemma}{Original lemma}{Original lemmas}
\makeatletter
\newlabel{algebra-lemma-quasi-compact}{{00E8}{source}{}{sourcepointer-2}{}}
\newlabel{algebra-lemma-quasi-compact@cref}{{[originallemma][0][]00E8}{[1][0][]source}{}{}{}}
\makeatother
\crefname{originallemma}{original lemma}{original lemmas}
\Crefname{originallemma}{Original lemma}{Original lemmas}
\makeatletter
\newlabel{lemma-cut-complex-in-two}{{19867}{source}{}{sourcepointer-3}{}}
\newlabel{lemma-cut-complex-in-two@cref}{{[originallemma][0][]19867}{[1][0][]source}{}{}{}}
\makeatother
\crefname{originallemma}{original lemma}{original lemmas}
\Crefname{originallemma}{Original lemma}{Original lemmas}
\makeatletter
\newlabel{lemma-finite-cohomology}{{15965}{source}{}{sourcepointer-4}{}}
\newlabel{lemma-finite-cohomology@cref}{{[originallemma][0][]15965}{[1][0][]source}{}{}{}}
\makeatother
\crefname{originallemma}{original lemma}{original lemmas}
\Crefname{originallemma}{Original lemma}{Original lemmas}
\makeatletter
\newlabel{lemma-glue-perfect}{{19077}{source}{}{sourcepointer-5}{}}
\newlabel{lemma-glue-perfect@cref}{{[originallemma][0][]19077}{[1][0][]source}{}{}{}}
\makeatother
\crefname{originallemma}{original lemma}{original lemmas}
\Crefname{originallemma}{Original lemma}{Original lemmas}
\makeatletter
\newlabel{lemma-lift-bounded-pseudo-coherent-to-perfect}{{20000}{source}{}{sourcepointer-6}{}}
\newlabel{lemma-lift-bounded-pseudo-coherent-to-perfect@cref}{{[originallemma][0][]20000}{[1][0][]source}{}{}{}}
\makeatother
\crefname{originallemma}{original lemma}{original lemmas}
\Crefname{originallemma}{Original lemma}{Original lemmas}
\makeatletter
\newlabel{lemma-pull-perfect}{{19017}{source}{}{sourcepointer-7}{}}
\newlabel{lemma-pull-perfect@cref}{{[originallemma][0][]19017}{[1][0][]source}{}{}{}}
\makeatother
\crefname{originallemma}{original lemma}{original lemmas}
\Crefname{originallemma}{Original lemma}{Original lemmas}
\makeatletter
\newlabel{lemma-summands-pseudo-coherent}{{16196}{source}{}{sourcepointer-8}{}}
\newlabel{lemma-summands-pseudo-coherent@cref}{{[originallemma][0][]16196}{[1][0][]source}{}{}{}}
\makeatother
\title{Perfection of a pseudo-coherent complex can be checked on localizations at primes or at maximal ideals}
\author{The Stacks Project Authors}
\date{stacks-project:more-algebra.tex:local-snapshot-20261012; GNU FDL 1.2 or later}
\begin{document}
\maketitle
\setcounter{subsection}{1}
\noindent\textit{Editorial note.} All delivered mathematical text is the author's own, copied from the pinned Stacks Project chapter file more-algebra.tex; no statement or proof step was rewritten, reordered or removed.\par\medskip
\noindent\textit{Editorial note.} The chapter file pinned-original-source.tex carries only \textbackslash{}input\{preamble\} on its first line, so no preamble text lives inside the chapter. The retained file sources/stacks-more-algebra/retained-preamble.tex is a mechanical, line-for-line copy of the project's single shared preamble as staged in the snapshot; nothing was added, removed or redefined and no package was added.\par\medskip
\noindent\textit{Editorial note.} The source URL is the source-level Stacks Project URL, as required when more than one card is produced from one source; the per-claim tag is carried by claim\_key.\par\medskip
\noindent\textit{Editorial note.} The author's displayed chain of implications contains a typographical repetition, reading '(1) => (2) <=> (3) <=> (3) <=> (4) <=> (5)' with (3) written twice; the display is reproduced unchanged rather than silently corrected. The author's typo 'we chooses' in lemma-hom-from-Af is not in this span. The ten references inside the counted proof span are recorded as non-load-bearing navigation pointers, and the card S15-alg-cut-complex-in-two independently counts the cutting lemma at its own lines.\par\medskip
\noindent\textit{Allowed standard prerequisites (editorial).}\begin{itemize}
\item The pullback-perfect lemma: perfection is preserved by base change and can be pulled back. invoked repeatedly at proof lines 20089-20096. Gives the localization implications (1) => (2), (3) => (2), (2) => (4) and (3) => (5).
\item The cutting lemma for a complex vanishing at a prime, and tor-amplitude of the upper truncation. invoked at proof lines 20105-20109. Produces the direct sum decomposition whose lower half must be shown to vanish.
\item Pseudo-coherence giving finite cohomology, and Nakayama's lemma. used at proof lines 20112-20116. Is the pair that makes the contradiction argument close.
\end{itemize}\medskip
\ifdefined\hypertarget\hypertarget{sourcepointer-1}{}\fi
\noindent\textit{Source reference.} Original lemma 00DV: focused-seeds/AGP/sources/stacks-snapshot-20261011/tags.txt, source lines 512--512. Cross-chapter navigation pointer into the Stacks Project; the standard fact it supplies is applied inside the counted proof span and the argument resting on it stays there.\par\medskip
\ifdefined\hypertarget\hypertarget{sourcepointer-2}{}\fi
\noindent\textit{Source reference.} Original lemma 00E8: focused-seeds/AGP/sources/stacks-snapshot-20261011/tags.txt, source lines 525--525. Cross-chapter navigation pointer into the Stacks Project; the standard fact it supplies is applied inside the counted proof span and the argument resting on it stays there.\par\medskip
\ifdefined\hypertarget\hypertarget{sourcepointer-3}{}\fi
\noindent\textit{Source reference.} Original lemma 19867: sources/stacks-more-algebra/pinned-original-source.tex, source lines 19867--19867. Same-chapter navigation pointer; the statement it names is standard for the section and the argument that consumes it stays inside the counted proof span.\par\medskip
\ifdefined\hypertarget\hypertarget{sourcepointer-4}{}\fi
\noindent\textit{Source reference.} Original lemma 15965: sources/stacks-more-algebra/pinned-original-source.tex, source lines 15965--15965. Same-chapter navigation pointer; the statement it names is standard for the section and the argument that consumes it stays inside the counted proof span.\par\medskip
\ifdefined\hypertarget\hypertarget{sourcepointer-5}{}\fi
\noindent\textit{Source reference.} Original lemma 19077: sources/stacks-more-algebra/pinned-original-source.tex, source lines 19077--19077. Same-chapter navigation pointer; the statement it names is standard for the section and the argument that consumes it stays inside the counted proof span.\par\medskip
\ifdefined\hypertarget\hypertarget{sourcepointer-6}{}\fi
\noindent\textit{Source reference.} Original lemma 20000: sources/stacks-more-algebra/pinned-original-source.tex, source lines 20000--20000. Same-chapter navigation pointer; the statement it names is standard for the section and the argument that consumes it stays inside the counted proof span.\par\medskip
\ifdefined\hypertarget\hypertarget{sourcepointer-7}{}\fi
\noindent\textit{Source reference.} Original lemma 19017: sources/stacks-more-algebra/pinned-original-source.tex, source lines 19017--19017. Same-chapter navigation pointer; the statement it names is standard for the section and the argument that consumes it stays inside the counted proof span.\par\medskip
\ifdefined\hypertarget\hypertarget{sourcepointer-8}{}\fi
\noindent\textit{Source reference.} Original lemma 16196: sources/stacks-more-algebra/pinned-original-source.tex, source lines 16196--16196. Same-chapter navigation pointer; the statement it names is standard for the section and the argument that consumes it stays inside the counted proof span.\par\medskip
\begin{lemma}
\label{lemma-check-perfect-stalks}
Let $R$ be a ring. Let $K^\bullet$ be a pseudo-coherent
complex of $R$-modules. Consider the following conditions
\begin{enumerate}
\item $K^\bullet$ is perfect,
\item for every prime ideal $\mathfrak p$ the complex
$K^\bullet \otimes_R R_{\mathfrak p}$ is perfect,
\item for every maximal ideal $\mathfrak m$ the complex
$K^\bullet \otimes_R R_{\mathfrak m}$ is perfect,
\item for every prime $\mathfrak p$ we have
$H^i(K^\bullet \otimes_R^{\mathbf{L}} \kappa(\mathfrak p)) = 0$ for all
$i \ll 0$,
\item for every maximal ideal $\mathfrak m$ we have
$H^i(K^\bullet \otimes_R^{\mathbf{L}} \kappa(\mathfrak m)) = 0$ for all
$i \ll 0$.
\end{enumerate}
We always have the implications
$$
(1) \Rightarrow (2) \Leftrightarrow (3) \Leftrightarrow (3)
\Leftrightarrow (4) \Leftrightarrow (5)
$$
If $K^\bullet$ is bounded below, then all conditions are equivalent.
\end{lemma}

\begin{proof}
By Lemma \ref{lemma-pull-perfect} we see that (1) implies (2).
It is immediate that (2) $\Rightarrow$ (3). Since every prime
$\mathfrak p$ is contained in a maximal ideal $\mathfrak m$,
we can apply Lemma \ref{lemma-pull-perfect} to the map
$R_\mathfrak m \to R_\mathfrak p$ to see that (3) implies (2).
Applying Lemma \ref{lemma-pull-perfect} to the residue maps
$R_\mathfrak p \to \kappa(\mathfrak p)$ and
$R_\mathfrak m \to \kappa(\mathfrak m)$ we see that (2) implies
(4) and (3) implies (5).

\medskip\noindent
Assume $R$ is local with maximal ideal $\mathfrak m$ and
residue field $\kappa$. We will show that if
$H^i(K^\bullet \otimes^\mathbf{L} \kappa) = 0$ for $i < a$
for some $a$, then $K$ is perfect. This will show that
(4) implies (2) and (5) implies (3) whence the first part
of the lemma. First we apply Lemma \ref{lemma-cut-complex-in-two}
with $i = a - 1$ to see that
$K^\bullet = \tau_{\leq a - 1}K^\bullet \oplus \tau_{\geq a}K^\bullet$
in $D(R)$ with $\tau_{\geq a}K^\bullet$ perfect of tor-amplitude
contained in $[a, \infty]$. To finish we need to show that
$\tau_{\leq a - 1}K$ is zero, i.e., that its cohomology groups are zero.
If not let $i$ be the largest index such that $M = H^i(\tau_{\leq a - 1}K)$
is not zero. Then $M$ is a finite $R$-module because
$\tau_{\leq a - 1}K^\bullet$ is pseudo-coherent
(Lemmas \ref{lemma-finite-cohomology} and \ref{lemma-summands-pseudo-coherent}).
Thus by Nakayama's lemma (Algebra, Lemma \ref{algebra-lemma-NAK})
we find that $M \otimes_R \kappa$ is nonzero.
This implies that
$$
H^i((\tau_{\leq a - 1}K^\bullet) \otimes_R^\mathbf{L} \kappa) =
H^i(K^\bullet \otimes_R^\mathbf{L} \kappa)
$$
is nonzero which is a contradiction.

\medskip\noindent
Assume the equivalent conditions (2) -- (5) hold and that
$K^\bullet$ is bounded below. Say $H^i(K^\bullet) = 0$ for
$i < a$. Pick a maximal ideal $\mathfrak m$ of $R$.
It suffices to show there exists an $f \in R$, $f \not \in \mathfrak m$
such that $K^\bullet \otimes_R^\mathbf{L} R_f$ is perfect
(Lemma \ref{lemma-glue-perfect} and
Algebra, Lemma \ref{algebra-lemma-quasi-compact}).
This follows from Lemma \ref{lemma-lift-bounded-pseudo-coherent-to-perfect}.
\end{proof}

\end{document}
